%% Template et conseils pour préparer des slides de présentation, pour
%% les étudiant·es de Licence Informatique et vidéoludisme de
%% l'université Paris 8.
%%
%% Par Farès Belhadj, amsi@up8.edu
%% Le 26 mai 2025
%% Personnalisé pour le projet tutoré "Exploration du Langage et du Cerveau"

\documentclass[t, aspectratio=169]{beamer}
\usepackage[utf8]{inputenc}
\usepackage[french]{babel}

%% Configuration de l'auteur pour le pied de page
\usetheme[auteur={N.~Roy \& T.~Massenya}, auteurbis={Nicolas Roy \& Théo Massenya}]{m2ise}

% Enable cancelto in math
\usepackage{cancel}
\renewcommand{\CancelColor}{\color{utred}}

% Funny icons et autres
% \usepackage[tikz]{bclogo}
\usepackage{listings}
\usepackage{float}
\usepackage{url}
\usepackage{datetime}
\usepackage{lmodern}
\usepackage{hyperref}
\usepackage{keystroke}
\usepackage{xcolor}
\usepackage{fancyvrb}
\usepackage{relsize}
\usepackage{ragged2e}
\usepackage{inconsolata}
\usepackage{amsmath, amssymb}

\lstset{
  language=Python,
  keywordstyle=\color{teal}\textbf,
  stringstyle=\color{blue},
  identifierstyle=\itshape,
  commentstyle=\itshape\color{gray},
  basicstyle=\ttfamily\tiny,
  numbers=none,
  frame=none,
  tabsize=2,
  breaklines=true,
  breakatwhitespace=true,
  inputencoding=utf8,
  extendedchars=true,
  literate={á}{{\'a}}1 {ã}{{\~a}}1 {à}{{\`a}}1 {é}{{\'e}}1 {ê}{{\^e}}1 {è}{{\`e}}1 {î}{{\^i}}1 {ç}{{\c c}}1 {ù}{{\`u}}1 {ô}{{\^o}}1 {°}{{\#}}1,
}

\newcommand{\gl}{\guillemotleft}
\newcommand{\gr}{\guillemotright}
\newcommand{\eg}[1]{~\gl~#1~\gr~}
\newcommand{\todo}{\texttt{TODO}}

\title{Exploration du Langage et du Cerveau via l'IA}
\subtitle{Rapport de Projet Tutoré}
\author{Nicolas Roy \& Théo Massenya}
\institute{Sous la direction de : Mme Revekka KYRIAKOGLOU\\Université Paris 8}
\date{Juin 2026}
\newcommand\code[1]{\textcolor{utred}{\texttt{#1}}}

\begin{document}

% Page de titre
\titleframe

% Table des matières / Plan
\begin{frame}{Plan de la présentation}
  \vspace{0.3cm}
  \begin{columns}[t]
    \begin{column}{0.48\textwidth}
      \begin{block}{1. Contexte et Objectifs}
        \scriptsize
        \begin{itemize}
          \item Le parallèle IA/Cerveau
          \item Problématique et hypothèses
        \end{itemize}
      \end{block}
      
      \begin{block}{2. État de l'Art}
        \scriptsize
        \begin{itemize}
          \item Modèle MUC (Hagoort)
          \item Dimensionnalité Intrinsèque (TwoNN)
          \item Le profil « Hunchback »
        \end{itemize}
      \end{block}
    \end{column}
    
    \begin{column}{0.48\textwidth}
      \begin{block}{3. Notre Approche et Méthodologie}
        \scriptsize
        \begin{itemize}
          \item Modèles et architecture
          \item Protocole Sémantique vs Jabberwocky
          \item Tokenisation (BPE) et alignement
        \end{itemize}
      \end{block}
      
      \begin{block}{4. Résultats et Analyse}
        \scriptsize
        \begin{itemize}
          \item Profils Hunchback et Ramping
          \item Discussion neurolinguistique
        \end{itemize}
      \end{block}
      
      \begin{block}{5. Conclusion}
        \scriptsize
        \begin{itemize}
          \item Validation \& Perspectives
        \end{itemize}
      \end{block}
    \end{column}
  \end{columns}
\end{frame}

% ==============================================================================
\section{Contexte et Objectifs}
% ==============================================================================

\begin{frame}{1. Contexte : Le parallèle IA et Cerveau}
  \begin{columns}[c]
    \begin{column}{0.75\textwidth}
      \begin{block}{L'essor des Transformers}
        \small
        \begin{itemize}
          \item Succès phénoménal des LLMs sur les tâches linguistiques complexes.
          \item Fonctionnement interne souvent décrit comme une \textcolor{utred}{\textbf{boîte noire}}.
        \end{itemize}
      \end{block}
      
      \begin{block}{Le rapprochement Cognitif}
        \small
        \begin{itemize}
          \item Les LLMs et le cerveau humain traitent le même signal : le langage naturel.
          \item \textbf{Question clé} : L'IA extrait-elle et structure-t-elle les représentations du sens de manière géométriquement analogue au cerveau ?
        \end{itemize}
      \end{block}
    \end{column}

  \end{columns}
\end{frame}

\begin{frame}{1. Problématique et Hypothèses}
  \begin{block}{Problématique}
    \centering\large\color{utred}
    Comment caractériser la dynamique d'unification sémantique au sein des couches d'un Transformer, et dans quelle mesure reflète-t-elle les signatures neuronales du cerveau humain ?
  \end{block}
  \vspace{0.2cm}
  \begin{columns}[t]
    \begin{column}{0.31\textwidth}
      \begin{block}{Hypothèse 1 : Ramping}
        \scriptsize
        Accumulation continue de l'information sémantique le long de la phrase, traduisant une augmentation de la dimensionnalité géométrique locale.
      \end{block}
    \end{column}
    
    \begin{column}{0.31\textwidth}
      \begin{block}{Hypothèse 2 : Cohérence}
        \scriptsize
        Une phrase cohérente engendre un surcoût d'unification progressive (ID plus élevée) par rapport à une suite vide de sens (Jabberwocky).
      \end{block}
    \end{column}
    
    \begin{column}{0.31\textwidth}
      \begin{block}{Hypothèse 3 : Hunchback}
        \scriptsize
        Le profil « bossu » d'ID (expansion puis compression) sera perturbé par l'absence de cohérence sémantique (Jabberwocky).
      \end{block}
    \end{column}
  \end{columns}
\end{frame}

% ==============================================================================
\section{État de l'Art}
% ==============================================================================

\begin{frame}{2. Principes Neurolinguistiques : Le Modèle MUC}
  \begin{columns}[c]
    \begin{column}{0.52\textwidth}
      \small
      Le modèle \textbf{Memory, Unification, and Control (MUC)} de Peter Hagoort décrit le traitement du langage dans l'hémisphère gauche :
      \vspace{0.2cm}
      \begin{itemize}
        \item \textbf{Memory (Cortex temporal)} : Stockage et récupération du lexique (structures syntaxiques et sémantiques).
        \item \textbf{Unification (Gyrus frontal inférieur / Broca)} : Assemblage dynamique et intégration des briques en une structure sémantique globale cohérente.
        \item \textbf{Control (Cortex préfrontal)} : Gestion de l'attention et des objectifs communicatifs.
      \end{itemize}
    \end{column}
    
    \begin{column}{0.44\textwidth}
      \centering
      \includegraphics[width=\linewidth]{The-MUC-model-of-language-The-figure-displays-a-lateral-view-of-the-left-hemisphere-The.png}
      \tiny \\ \textit{Localisation anatomique des composantes du MUC.}
    \end{column}
  \end{columns}
\end{frame}

\begin{frame}{2. La Dimensionnalité Intrinsèque (ID) \& TwoNN}
  \begin{columns}[c]
    \begin{column}{0.5\textwidth}
      \begin{block}{L'Hypothèse de la Variété (Manifold)}
        \small
        Bien que représentées dans des espaces de très haute dimension (ex: $D = 768$ pour GPT-2), les données réelles sont confinées sur des structures locales de dimension beaucoup plus faible : l'\textbf{ID}.
      \end{block}
      
      \begin{block}{L'estimateur TwoNN (Facco et al., 2017)}
        \small
        \begin{itemize}
          \item Calcule le ratio des distances aux deux plus proches voisins pour chaque point : $\mu_i = r_{i,2} / r_{i,1}$.
          \item ID estimée à partir de la fonction de répartition empirique de $\mu$.
          \item Indépendant des variations de densité locale.
        \end{itemize}
      \end{block}
    \end{column}
    
    \begin{column}{0.46\textwidth}
      \begin{exampleblock}{Propriété géométrique clé}
        \scriptsize
        L'apprentissage prédictif (minimiser l'erreur du prochain token) force la compression géométrique des représentations sur des variétés compactes et de basse dimensionnalité intrinsèque.
      \end{exampleblock}
      \centering
      \vspace{0.5cm}
      \large $d = - \frac{\ln(1 - F(\mu))}{\ln(\mu)}$
    \end{column}
  \end{columns}
\end{frame}

\begin{frame}{2. Dynamique par couches : Le Profil « Hunchback »}
  \begin{columns}[c]
    \begin{column}{0.5\textwidth}
      \small
      Ansuini et al. (2019) mettent en évidence une évolution non-linéaire de l'ID à travers les couches des réseaux profonds :
      \vspace{0.2cm}
      \begin{itemize}
        \item \textbf{1. Phase d'expansion} : Dans les premières couches, l'ID augmente pour « déplier » les données d'entrée et séparer les facteurs de variation.
        \item \textbf{2. Phase de compression} : Dans les couches supérieures, l'ID diminue drastiquement pour ne retenir que l'information utile à la prédiction (goulot d'étranglement sémantique).
      \end{itemize}
    \end{column}
    
    \begin{column}{0.46\textwidth}
      \centering
      \includegraphics[width=0.9\linewidth]{profil Hunchback.png}
      \tiny \\ \textit{Profil Hunchback typique dans VGG-16 (Ansuini et al.).}
    \end{column}
  \end{columns}
\end{frame}

\begin{frame}{2. Signatures Temporelles de l'Unification}
  \begin{block}{Le phénomène de « Ramping »}
    \small
    \begin{itemize}
      \item Chez l'humain (EEG/fMRI), l'unification s'accompagne d'une augmentation graduelle et continue de l'activité neuronale dans l'aire de Broca au cours de la phrase.
      \item Desbordes et al. (2023) montrent que les représentations internes des LLMs évoluent de manière similaire : l'ID augmente progressivement mot après mot.
      \item \textbf{Stimuli sémantiques vs non-sensés} : Le Ramping est soutenu par la cohérence sémantique. L'absence de sens (ex: Jabberwocky) perturbe cette accumulation progressive.
    \end{itemize}
  \end{block}
  \begin{center}
    \textcolor{utred}{\textbf{Hypothèse}} : Le ramping géométrique est une signature de l'effort d'intégration sémantique.
  \end{center}
\end{frame}

% ==============================================================================
\section{Notre Approche - Méthodologie et Protocole}
% ==============================================================================

\begin{frame}{3. Outils et Modèle Expérimental}
  \begin{columns}[t]
    \begin{column}{0.48\textwidth}
      \begin{block}{Modèle de Langage : GPT-2 Small}
        \small
        \begin{itemize}
          \item Modèle \textit{Decoder-only} causal (12 couches, $D = 768$).
          \item \textbf{Justification} : Le masque causal traite le texte de gauche à droite, simulant la linéarité temporelle de la lecture humaine.
        \end{itemize}
      \end{block}
    \end{column}
    
    \begin{column}{0.48\textwidth}
      \begin{block}{Framework Technique}
        \small
        \begin{itemize}
          \item \textbf{PyTorch} et \textbf{Hugging Face} pour l'extraction fine des \textit{hidden states} par couche et par token.
          \item \textbf{scikit-dimension} pour l'implémentation robuste de TwoNN.
          \item \textbf{spaCy} pour l'analyse syntaxique (POS Tagging) et la génération.
        \end{itemize}
      \end{block}
    \end{column}
  \end{columns}
\end{frame}

\begin{frame}{3. Protocole Expérimental : Sémantique vs Jabberwocky}
  \begin{block}{Objectif}
    Isoler l'impact du traitement sémantique en maintenant la structure grammaticale et le patron syntaxique strictement identiques.
  \end{block}
  \vspace{0.1cm}
  \begin{columns}[t]
    \begin{column}{0.48\textwidth}
      \begin{block}{Groupe Sémantique (N = 500)}
        \scriptsize
        \begin{itemize}
          \item Phrases extraites de \textbf{WikiText-2} (Wikipedia).
          \item Langue : \textbf{Anglais} (GPT-2 Small étant entraîné sur de l'anglais).
          \item Longueur : 8 à 20 mots, syntaxe propre.
        \end{itemize}
      \end{block}
    \end{column}
    
    \begin{column}{0.48\textwidth}
      \begin{block}{Groupe Jabberwocky (N = 500)}
        \scriptsize
        \begin{itemize}
          \item Remplacement systématique des \textbf{mots porteurs} (noms, verbes, adjectifs, adverbes) par des pseudo-mots respectant la phonotactique anglaise.
          \item Conservation stricte des \textbf{mots-outils} (prépositions, articles, pronoms) et de la syntaxe.
        \end{itemize}
      \end{block}
    \end{column}
  \end{columns}
  \vspace{0.2cm}
  \centering
  \fbox{\small \textbf{Exemple} : \textit{"The evolution of language is taught"} $\rightarrow$ \textit{"The zumplation of language is carting"}}
\end{frame}

\begin{frame}{3. Enjeux de la Tokenisation : La Fragmentation BPE}
  \begin{columns}[c]
    \begin{column}{0.52\textwidth}
      \begin{block}{Le biais de la tokenisation}
        \small
        \begin{itemize}
          \item L'algorithme BPE fragmente les mots inconnus (pseudo-mots) en sous-tokens (ex: \textit{zumplation} $\rightarrow$ \textit{zump} + \textit{lation}).
          \item Le Jabberwocky comporte en moyenne \textcolor{utred}{\textbf{2,05 fois plus de tokens}} que le sémantique.
        \end{itemize}
      \end{block}
      
      \begin{block}{Solution : Alignement Temporel Relative}
        \small
        \begin{itemize}
          \item Normalisation de la longueur par extraction d'un token représentatif par \textbf{décile de position} (10\% à 100\%).
          \item Neutralise le biais du nombre de tokens sur le calcul du ramping.
        \end{itemize}
      \end{block}
    \end{column}
    
    \begin{column}{0.44\textwidth}
      \centering
      \includegraphics[width=\linewidth]{fig3_bpe_fragmentation.png}
      \tiny \\ \textit{Histogramme du ratio de fragmentation BPE.}
    \end{column}
  \end{columns}
\end{frame}

\begin{frame}{3. Expérience Classe Grammaticale}

\end{frame}

\begin{frame}{4. Pipeline d'Implémentation du Projet}
  \centering
  \begin{tikzpicture}[node distance=1.5cm, every node/.style={fill=white, draw=utred, rounded corners, font=\scriptsize, align=center}, arrow/.style={thick, ->, >=stealth, color=utblack}]
    % Nodes
    \node (v2) [draw=utblack] {0. Validation TwoNN \\ (Données synthétiques)};
    \node (s1) [below of=v2, yshift=0.3cm] {1. Génération Stimuli \\ (Normalisation \& spaCy)};
    \node (s2) [below of=s1, yshift=0.3cm] {2. Extraction Hidden States \\ (GPT-2 Small + Alignement Déciles)};
    \node (s3) [below of=s2, yshift=0.3cm] {3. Calcul de l'ID locale \\ (Algorithme TwoNN + Bootstrap)};
    \node (s4) [below of=s3, yshift=0.3cm, fill=utred!10] {4. Visualisation \& Tests Statistiques \\ (Mann-Whitney / Regressions)};

    % Arrows
    \draw [arrow] (v2) -- (s1);
    \draw [arrow] (s1) -- (s2);
    \draw [arrow] (s2) -- (s3);
    \draw [arrow] (s3) -- (s4);
  \end{tikzpicture}
\end{frame}

% ==============================================================================
\section{Résultats et Analyse}
% ==============================================================================


\begin{frame}{4. Évolution par Couche : La Rupture du Hunchback}
  \begin{columns}[c]
    \begin{column}{0.5\textwidth}
      \begin{block}{Sémantique (Bleu) : Hunchback Typique}
        \scriptsize
        \begin{itemize}
          \item \textbf{Expansion} : L'ID augmente de 18.65 à \textbf{21.14} (couche 6). Démêlage des variables sémantiques.
          \item \textbf{Compression} : L'ID diminue jusqu'à \textbf{18.29} (couche 12). Extraction d'un concept abstrait.
        \end{itemize}
      \end{block}
      
      \begin{block}{Jabberwocky (Rouge) : Explosion Finale}
        \scriptsize
        \begin{itemize}
          \item L'ID ne compresse jamais : elle croît continuellement pour culminer à \textbf{30.01} à la couche 12.
          \item Échec de l'unification : faute de cohérence sémantique, la représentation reste bruitée.
        \end{itemize}
      \end{block}
    \end{column}
    
    \begin{column}{0.48\textwidth}
      \centering
      \includegraphics[width=\linewidth]{fig1_hunchback.png}
      \tiny \\ \textit{ID par couche sur le dernier token de la phrase.}
    \end{column}
  \end{columns}
\end{frame}

\begin{frame}{4. Dynamique Temporelle : Le Ramping de l'ID}
  \begin{columns}[c]
    \begin{column}{0.46\textwidth}
      \begin{block}{Sémantique : Ramping précoce}
        \scriptsize
        \begin{itemize}
          \item Pente positive et statistiquement significative dès la **couche 3** et jusqu'à la couche 11.
          \item \textbf{Processus d'intégration continue} : l'unification sémantique augmente progressivement la charge représentationnelle.
        \end{itemize}
      \end{block}
      
      \begin{block}{Jabberwocky : Profil en 2 étapes}
        \scriptsize
        \begin{itemize}
          \item \textbf{Holding Pattern} (10\% à 80\%) : ID basse et stable. Le modèle ne stocke que la structure syntaxique (peu coûteuse).
          \item \textbf{Wrap-up final} (90\%-100\%) : Bond d'ID brutal. Au point final, la fermeture de la proposition force l'unification, provoquant l'explosion dimensionnelle.
        \end{itemize}
      \end{block}
    \end{column}
    
    \begin{column}{0.52\textwidth}
      \centering
      \includegraphics[width=\linewidth]{fig2_ramping.png}
      \tiny \\ \textit{Évolution temporelle de l'ID par déciles de position.}
    \end{column}
  \end{columns}
\end{frame}

\begin{frame}{4. Confrontation Cognitive : Le rôle du Contrôle}
  \begin{block}{Le Parallèle MUC}
    \small
    Le Ramping sémantique progressif et l'augmentation géométrique dans les couches intermédiaires modélisent de manière convaincante le processus d'\textbf{Unification} de l'aire de Broca.
  \end{block}
  \vspace{0.1cm}
  \begin{columns}[t]
    \begin{column}{0.48\textwidth}
      \begin{block}{Le Cerveau Humain}
        \scriptsize
        \begin{itemize}
          \item Face au Jabberwocky, le cerveau humain désengage rapidement son attention sémantique après constatation de l'incohérence.
          \item Ce rôle est joué par le **Contrôle** cognitif (Cortex préfrontal dorsolatéral).
        \end{itemize}
      \end{block}
    \end{column}
    
    \begin{column}{0.48\textwidth}
      \begin{block}{Le Transformer (GPT-2)}
        \scriptsize
        \begin{itemize}
          \item Dépourvu de régulation attentionnelle cognitive supérieure.
          \item Forcé d'appliquer le même mécanisme uniforme, il s'épuise à chercher un sens aux pseudo-mots, d'où l'explosion de l'ID.
        \end{itemize}
      \end{block}
    \end{column}
  \end{columns}
\end{frame}

% ==============================================================================
\section{Conclusion}
% ==============================================================================

\begin{frame}{5. Conclusion : Validation des Hypothèses}
  \begin{block}{Validation des hypothèses de travail}
    \small
    \begin{enumerate}
      \item \textbf{Ramping validé} : La dimensionnalité géométrique augmente de manière lente et continue au cours de la lecture sémantique.
      \item \textbf{Effet de cohérence validé} : L'unification sémantique présente une charge géométrique progressive, confirmée par la comparaison avec le Jabberwocky.
      \item \textbf{Déformation Hunchback validée} : L'absence de sens bloque la compression géométrique dans les couches finales cachées de GPT-2 Small.
    \end{enumerate}
  \end{block}
  \vspace{0.2cm}
  \begin{center}
    \small \textbf{Contribution} : L'ID se révèle être un excellent marqueur géométrique macroscopique pour sonder le traitement sémantique des modèles de langage.
  \end{center}
\end{frame}

\begin{frame}{5. Limites et Perspectives}
  \begin{columns}[t]
    \begin{column}{0.48\textwidth}
      \begin{block}{Limites de l'étude}
        \small
        \begin{itemize}
          \item Modèle restreint en taille (GPT-2 Small, 117M params).
          \item Restreint à la langue anglaise en raison de l'entraînement originel du modèle.
        \end{itemize}
      \end{block}
    \end{column}
    
    \begin{column}{0.48\textwidth}
      \begin{block}{Perspectives de recherche}
        \small
        \begin{itemize}
          \item \textbf{Changement d'échelle} : Reproduire le protocole sur des modèles plus volumineux (ex: Llama 3.1 8B).
          \item \textbf{Estimateurs d'ID} : Tester d'autres algorithmes locaux (ex: Levina-Bickel, FisherS) pour consolider les mesures géométriques.
        \end{itemize}
      \end{block}
    \end{column}
  \end{columns}
\end{frame}

\lastframe%
\end{document}
